\section{Introduction}
\label{sec:intro}

Fraud on online platforms---fake reviews, money laundering, collusive crowdsourcing---causes billions in annual losses. Detecting these activities is inherently a graph problem: fraudsters interact with victims, form collusive rings, and leave topological traces invisible to non-relational models~\cite{wang2019fdgars,dou2020enhancing}. Graph Neural Networks (GNNs) have become the dominant approach because they jointly model node attributes and neighborhood structure~\cite{kipf2016semi,hamilton2017inductive}.

However, fraud detection presents two hard challenges for standard GNNs. First, camouflage: fraudsters mimic benign feature distributions to blend in~\cite{dou2020enhancing}. Second, neighborhood inconsistency: a fraudster's neighbors are mostly benign, diluting the fraud signal during message passing~\cite{liu2020alleviating}. On top of these, class imbalance is severe---fraud typically constitutes 3--17\% of nodes, so a model predicting ``benign'' for all achieves high accuracy while catching zero fraud. Existing detectors address these via neighbor filtering or sampling~\cite{dou2020enhancing,liu2021pick,liu2020alleviating}, adding complexity without treating pattern rarity itself. Focal loss~\cite{lin2017focal} reweights by classification difficulty, not intrinsic rarity.

GAAP~\cite{duan2025gaap} departs from neighbor engineering by learning \textit{attribute-association patterns}: adaptive feature binning preserves per-attribute semantics, message passing extracts neighborhood association context, and a global cross-attention module captures long-range correlations. GAAP achieves SOTA against 24 baselines on 7 datasets without neighbor filtering. However, it trains with standard cross-entropy, assigning equal weight to every node. Common benign patterns dominate the gradient, while rare fraud-indicative patterns receive negligible attention. Our central question: can we improve GAAP simply by telling the optimizer which training examples encode rare patterns?

We propose RP-GAAP (Rare-Pattern-weighted GAAP), a lightweight extension that computes per-sample weights from quantile-binned feature pattern frequencies and uses these scores in the cross-entropy loss. The method adds three hyperparameters, introduces no learned parameters, and requires no architectural changes. We evaluate RP-GAAP on four fraud detection benchmarks (YelpChi, T-Finance, Elliptic, Tolokers) across four loss configurations.
